\documentclass{article}
\usepackage{graphicx} % Required for inserting images
\usepackage{hyperref} % For professional links and table of contents
\hypersetup{colorlinks=true, linkcolor=blue, urlcolor=blue}

\title{Economy and social Change}
\author{Bisan Jabareen}
\date{Spring 2026}
\begin{document}


\maketitle
\newpage
\tableofcontents
\newpage
\section{Lecture 1: Introduction to Economic Anthropology}

\subsection{Etymology and Definition}
The word \textbf{economy} originates from Greece (\textit{oikonomia}). Historically, the term was focused on household affairs.

\subsection{Economics vs. Economic Anthropology}
\begin{itemize}
    \item \textbf{Economists:} Emphasize what money \textit{does}—a medium of exchange, reserve fund, or means of accounting.
    \item \textbf{Anthropologists:} View the economy as an integral part of the hierarchies and networks of exchange through which societies circulate. You cannot separate economy from the society in which it operates (embeddedness).
\end{itemize}

\textit{Note: Economist and Nobel Prize winner Richard Thaler has noted that people often do not make rational decisions regarding money.}

\subsection{Diversity vs. Universal}
\begin{itemize}
    \item \textbf{Livelihoods:} Humans have the same basic needs, but the way we meet these means is what is studied and what differs.
    \item \textbf{Focus:} Anthropologists focus on how people \textbf{produce, exchange, and consume} material objects.
\end{itemize}

\subsection{Economic anthropology - Some History}
\begin{itemize}
    Karl Polanyi argues that the economy is embedded in society. The economy is not a separate sphere; it is part of the social fabric, but that Capitalism is a disembedded system.
\end{itemize}

\subsection{Capitalism and Society}
\begin{itemize}
    \item \textbf{Time/Productivity:} Rest in capitalism is often viewed as laziness; there is a constant pressure to ``do something with your time.''
    \item \textbf{Embedded vs. Disembedded:} 
    \begin{itemize}
        \item \textbf{Embedded:} A grandma making soup when you are sick (moral obligation).
        \item \textbf{Disembedded:} A restaurant giving you soup only if you have money (self-interest).
    \end{itemize}
\end{itemize}

\subsection{Two Main Approaches to Financial Decisions}
\begin{enumerate}
    \item \textbf{Formalists:} We have the choice and make decisions based on personal motives and self-interest. 
    \textit{(Example: Buying an iPhone 17 because you personally wanted to.)}
    \item \textbf{Substantivists:} Decisions are funneled by society; you are a victim of capitalism. 
    \textit{(Example: Buying an iPhone 17 because society pressures you to.)}
\end{enumerate}

\subsection{Three Phases of Economic Activity}
Production $\rightarrow$ Exchange $\rightarrow$ Consumption

\subsection{Karl Marx}
Marx argued that work is our most basic human activity. Human consciousness is determined by work (\textbf{labor}).
\begin{itemize}
    \item \textbf{Mode of Production:} Introducted by Karl Marx and Friedrich Engels to
describe the way societies organize production,
combining forces of production (tools, labor, technology)
with relations of production (social and property
relations).
\end{itemize}

\subsection{Eric Wolf: Three Modes of Production}
\begin{enumerate}
    \item \textbf{Domestic (Kin-ordered):} The whole chain of production is in the family (e.g., subsistence farming).
    \item \textbf{Tributary:} Paying tribute to rulers; mandatory giving of goods. Today, this manifests as taxes paid to the government.
    \item \textbf{Capitalist:} Humans do \textbf{not} own their means of production. People sell their labor for money. Implications include a lack of control over work conditions (WLB, vacation days, work hours) and business decisions (budgets, projects, etc). This is considered a disembedded system, according to Karl Polanyi.
\end{enumerate}

\end{document}
